
\documentclass{article}
\usepackage{colortbl}
\usepackage{makecell}
\usepackage{multirow}
\usepackage{supertabular}

\begin{document}

\newcounter{utterance}

\centering \large Interaction Transcript for game `taboo', experiment `high\_en', episode 9 with reflexion{-}Qwen3.5{-}4B.
\vspace{24pt}

{ \footnotesize  \setcounter{utterance}{1}
\setlength{\tabcolsep}{0pt}
\begin{supertabular}{c@{$\;$}|p{.15\linewidth}@{}p{.15\linewidth}p{.15\linewidth}p{.15\linewidth}p{.15\linewidth}p{.15\linewidth}}
    \# & $\;$A & \multicolumn{4}{c}{Game Master} & $\;\:$B\\
    \hline

    \theutterance \stepcounter{utterance}  
    & & \multicolumn{4}{p{0.6\linewidth}}{
        \cellcolor[rgb]{0.9,0.9,0.9}{
            \makecell[{{p{\linewidth}}}]{
                \texttt{\tiny{[P1$\langle$GM]}}
                \texttt{You have attempted this same task before and did not succeed. These are the plans you wrote after those attempts.} \\
\texttt{Trial \#1: New plan: I failed because I provided a definition ("The opposite of 'most' or 'many'") rather than a direct clue that forces the other player to guess the specific word "least" within three tries. The target word "least" is a superlative form of "little," and the related word "little" is explicitly provided. I should have used the related word "little" as the primary anchor in my clue, perhaps phrasing it as "The superlative form of 'little' meaning 'not much'." This leverages the provided related word to guide the guesser directly to "least" instead of just describing the concept.} \\
\\ 
\texttt{Trial \#2: New plan: I failed because my clue was too abstract and relied on a definition rather than a direct linguistic connection to the provided related words. The target word "least" is explicitly the superlative form of the related word "little." I should have used the related word "little" as the primary anchor in my description, specifically stating "The superlative form of 'little' meaning the smallest amount." This directly leverages the allowed related word to force the guesser to recognize "least" immediately, rather than just describing the concept of a minimum quantity.} \\
\\ 
\\ 
\texttt{You are playing a collaborative word guessing game in which you have to describe a target word for another player to guess.} \\
\\ 
\texttt{Rules:} \\
\texttt{(a) You have to reply in the form: CLUE: <some text>. Guesses from the other player will start with GUESS.} \\
\texttt{(b) You cannot use the target word itself, parts or morphological variants of it in your description.} \\
\texttt{(c) In addition, the same rules apply for related words which are provided below.} \\
\\ 
\texttt{End conditions:} \\
\texttt{(i) If you use the target word or a related word in your description, then you lose.} \\
\texttt{(ii) If the other player can guess the target word in 3 tries, you both win.} \\
\\ 
\texttt{Let us start.} \\
\\ 
\texttt{This is the target word that you need to describe and that the other player needs to guess:} \\
\\ 
\texttt{least} \\
\\ 
\texttt{Related words are:} \\
\\ 
\texttt{{-} amount} \\
\texttt{{-} most} \\
\texttt{{-} little} \\
\\ 
\texttt{Important: You are under time pressure, give short descriptions that are to the point!} \\
            }
        }
    }
    & \\ \\

    \theutterance \stepcounter{utterance}  
    & \multicolumn{4}{p{0.6\linewidth}}{
        \cellcolor[rgb]{0.9,0.9,0.9}{
            \makecell[{{p{\linewidth}}}]{
                \texttt{\tiny{[P1$\rangle$GM]}}
                \texttt{CLUE: The superlative form of "little" meaning the smallest amount.} \\
            }
        }
    }
    & & \\ \\

    \theutterance \stepcounter{utterance}  
    & & & \multicolumn{2}{p{0.3\linewidth}}{
        \cellcolor[rgb]{0.9,0.9,0.9}{
            \makecell[{{p{\linewidth}}}]{
                \texttt{\tiny{[GM$|$GM]}}
                \texttt{continue} \\
            }
        }
    }
    & & \\ \\

    \theutterance \stepcounter{utterance}  
    & & & \multicolumn{2}{p{0.3\linewidth}}{
        \cellcolor[rgb]{0.9,0.9,0.9}{
            \makecell[{{p{\linewidth}}}]{
                \texttt{\tiny{[GM$|$GM]}}
                \texttt{Related word 'little' (stem=littl) is similar to clue word 'little' (stem=littl)} \\
            }
        }
    }
    & & \\ \\

    \theutterance \stepcounter{utterance}  
    & \multicolumn{4}{p{0.6\linewidth}}{
        \cellcolor[rgb]{0.95,0.95,0.95}{
            \makecell[{{p{\linewidth}}}]{
                \texttt{\tiny{[P2$\rightarrow$P1]}}
                \texttt{Post{-}episode Reflexion (not sent to GM):} \\
\texttt{New plan: I failed because I provided a definition ("The superlative form of 'little' meaning the smallest amount") which, while factually correct, is too abstract and relies on the guesser's internal knowledge rather than explicitly linking to the provided related word "little" in a way that forces the specific target "least". The related word "little" is explicitly given, so I must use it as the direct anchor. I will phrase my clue to explicitly state that the target is the superlative form of the related word "little" (e.g., "The superlative form of 'little'"), ensuring the guesser connects the provided hint directly to the target word without needing to derive the definition from scratch.} \\
            }
        }
    }
    & & \\ \\

\end{supertabular}
}

\end{document}
